\section{The single-cluster bound in Hastings' cluster expansion}
\label{sec:hastings_single_cluster_bound}

\begin{theorem}[The single-cluster bound]
    \label{thm:hastings_single_cluster_bound}
    Let $H$ be a Hermitian Hamiltonian, let $\beta\ge0$, and put
    $Z=\operatorname{tr}e^{-\beta H}$. For a cluster $C$ and its disjoint boundary
    $B(C)$, suppose that the Hermitian terms $H_k$ have operator norm at most $J$.
    As in the source, the terms are indexed by sites, with $H_k$ supported near
    site $k$, so $C$ and $B(C)$ are sets of sites. Write $n_C=|C|$ and
    $n_{B(C)}=|B(C)|$ for their numbers of sites, $H_C=\sum_{k\in C}H_k$ and
    $H_{B(C)}=\sum_{k\in B(C)}H_k$. Then for every operator $O$
    \begin{align}
        Z^{-1}\left|\operatorname{tr}\left(Oe^{-\beta(H-H_C-H_{B(C)})}\right)\right|
        \le\|O\|e^{(n_C+n_{B(C)})J\beta}.
        \label{eq:hastings_single_cluster_application}
    \end{align}
    This is the single-cluster estimate of~\cite[Section IV]{Hastings2006Solving},
    with an absolute value on the trace.
\end{theorem}
\begin{proof}
    It is the finite-sum perturbation estimate in the companion quantum-channel
    volume~\cite[``The single-cluster bound'']{QICLean2026Blueprint}, applied to
    $T=C\cup B(C)$. The source proves it by an imaginary-time Dyson expansion. For
    the companion proof, put $V=H_C+H_{B(C)}$ and $a=(n_C+n_{B(C)})J$, and let $I$ be the identity.
    The triangle inequality gives $\|V\|\le a$. Hermiticity and monotonicity of the
    trace exponential then give
    \begin{align}
        V & \le aI,                                                 \\
        \operatorname{tr}e^{-\beta H+\beta V}
          & \le\operatorname{tr}e^{-\beta H+\beta aI}=e^{\beta a}Z.
        \label{eq:hastings_single_cluster_trace_bound}
    \end{align}
    The matrix $A=e^{-\beta(H-V)}$ satisfies $A\ge0$, so
    \begin{align}
        |\operatorname{tr}(OA)|
         & \le\|O\|\operatorname{tr}A
        \le\|O\|e^{\beta a}Z,
        \label{eq:hastings_single_cluster_trace_pairing}
    \end{align}
    where the second inequality is~(\ref{eq:hastings_single_cluster_trace_bound}).
    For positive dimension, $Z>0$, and dividing
    (\ref{eq:hastings_single_cluster_trace_pairing}) by $Z$ gives
    (\ref{eq:hastings_single_cluster_application}). In dimension zero the estimate
    uses $0^{-1}=0$.
\end{proof}

\begin{definition}[The cluster part]
    \label{def:hastings_cluster_part}
    \lean{TNLean.Lattice.ClusterExpansion.clusterPart,
        TNLean.Lattice.ClusterExpansion.clusterPart_eq_sum_powerset}
    \leanok
    For a Hamiltonian $H$ containing the terms $H_k$, a cluster $C$, its boundary
    $B(C)$ and $\beta\in\mathbb R$, put
    \begin{align}
        \rho(\beta,C)=\sum_{D\subseteq C}(-1)^{|D|}e^{-\beta(H-H_{B(C)}-H_D)}.
        \label{eq:hastings_cluster_part}
    \end{align}
    By inclusion and exclusion this is the sum of the terms of the power series
    of $e^{-\beta H}$ that contain every $H_i$ with $i\in C$ and no $H_j$ with
    $j\in B(C)$, the operator $\rho(\beta,C)$
    of~\cite[lines 1174--1177]{Hastings2006Solving}. With $E=C\setminus D$ it is
    $\sum_{E\subseteq C}(-1)^{|C|-|E|}e^{-(Y+\sum_{i\in E}\beta H_i)}$ for
    $Y=\beta(H-H_C-H_{B(C)})$.
\end{definition}

\begin{theorem}[The bound on the cluster part]
    \label{thm:hastings_cluster_part_bound}
    \lean{TNLean.Lattice.ClusterExpansion.norm_trace_mul_clusterPart_le}
    \leanok
    \uses{def:hastings_cluster_part}
    Let $H$ be Hermitian, let the terms $H_k$ with $k\in C\cup B(C)$ be Hermitian
    with $\|H_k\|\le J$, and let $\beta\ge0$. Then for every operator $O$,
    with $Z=\operatorname{tr}e^{-\beta H}$,
    \begin{align}
        \left|\operatorname{tr}\bigl(O\rho(\beta,C)\bigr)\right|
        \le\|O\|\bigl(e^{\beta J}-1\bigr)^{n_C}e^{(n_C+n_{B(C)})J\beta}Z.
        \label{eq:hastings_cluster_part_bound}
    \end{align}
    This is the bound stated before equation (percb2)
    in~\cite[lines 1174--1195]{Hastings2006Solving}, with $\|O\|Z$ in place of the
    printed $\operatorname{tr}(Oe^{-\beta H})$, which is the form used in (percb2),
    and with $\|H_C+H_{B(C)}\|\le(n_C+n_{B(C)})J$. The last step of (percb2),
    $n_C+n_{B(C)}\le Sn_C$, is a property of clusters on a lattice and is not
    made here.
\end{theorem}

\begin{proof}
    \leanok
    \uses{def:hastings_cluster_part}
    By~(\ref{eq:hastings_cluster_part}) with $E=C\setminus D$, the alternating-sum
    bound of the companion quantum-channel
    volume~\cite[``Alternating sums of shifted exponentials'']{QICLean2026Blueprint},
    applied with $Y=\beta(H-H_C-H_{B(C)})$ and the terms $\beta H_i$, $i\in C$, of
    norm at most $\beta J$, gives
    \begin{align*}
        \left|\operatorname{tr}\bigl(O\rho(\beta,C)\bigr)\right|
        \le\|O\|\bigl(e^{\beta J}-1\bigr)^{n_C}\operatorname{tr}e^{-Y}.
    \end{align*}
    Since $-Y=-\beta H+\beta(H_C+H_{B(C)})$ and
    $\|\beta(H_C+H_{B(C)})\|\le(n_C+n_{B(C)})J\beta$, the bound for a Hermitian
    perturbation in the exponent in the same
    volume~\cite[``A Hermitian perturbation in the exponent'']{QICLean2026Blueprint}
    gives $\operatorname{tr}e^{-Y}\le e^{(n_C+n_{B(C)})J\beta}Z$.
\end{proof}
